\section{Related Work}
\label{sec:related-work}

Classical UAV path planning includes graph search, incremental search,
sampling-based planning, and potential-field methods. Dijkstra provides optimal
paths for non-negative edge costs \cite{dijkstra1959note}. A* adds an admissible
heuristic and normally
expands fewer nodes while preserving optimality \cite{hart1968astar}. D* Lite
reuses earlier search information after edge-cost changes, making it a natural
reference for dynamic routing rather than treating repeated full search as the
only classical option \cite{koenig2005dstarlite}. Broader UAV-planning surveys
emphasize the continued value of completeness, interpretability, and predictable
behavior in classical planners \cite{ghambari2024survey,venu2025comprehensive}.

RL methods instead learn a policy through interaction. DQN scales value-based
learning with neural function approximation \cite{mnih2015human}; Double DQN
reduces max-operator overestimation \cite{vanhasselt2016double}; and HER
relabels failed trajectories with achieved goals to improve learning in sparse
goal-conditioned tasks \cite{andrychowicz2017her}. UAV studies have applied
value-based and deep RL to path planning and obstacle avoidance, reporting
adaptability but also dependence on representation, reward, and evaluation
conditions \cite{tu2023uav,guo2023uav,zhao2022path,rocha2024dynamic}. Recent
work also combines temporal observations, recurrent policy learning, and an
explicit safety shield for dynamic obstacle avoidance
\cite{liu2026temporal}. Those richer sensing and safety assumptions are
complementary to this study's narrower goal: isolating algorithmic behavior
under a controlled shared graph and change-event contract.

Evaluation practice is especially important for deep RL. Henderson et al.
showed that implementation and random-seed variation can materially alter
reported conclusions \cite{henderson2018deep}. A test distribution of many
routes from one trained network does not quantify this training variance.
Accordingly, this study treats independently trained seeds as the uncertainty
units for RL and retains scenario-level pairing within each seed.

Systematic comparisons likewise identify different trade-offs in
predictability, adaptation, and computational cost between classical and
learned approaches \cite{mannan2023classical}. Prior UAV comparisons often vary
the map, metrics, or hardware between methods,
which makes algorithmic attribution difficult. This work focuses on controlled
same-scenario comparison: all methods receive the same persisted graph and
change events, while the benchmark separately varies endpoint, layout,
density, dynamics, scale, energy cost, no-fly behavior, and observation noise.
The aim is not to claim that a grid benchmark captures complete UAV operation,
but to identify precisely where differences arise before moving to richer
flight models.
